\section{Adapting nanoGPT to MAESTRO}
\label{sec:implementation}
nanoGPT provides compact PyTorch model and training code \citep{karpathynanogpt}. It lets us inspect data loading, optimization, and checkpointing and vary model size without changing frameworks. We choose it for clarity and reuse, without claiming it is the fastest implementation.

The decoder learns token and absolute position embeddings. Causal attention uses earlier tokens while masking future targets \citep{vaswani2017attention}. Blocks use pre-layer normalization, residual connections, and GELU feed-forward networks with four times the hidden width. Token embeddings and the output head share weights. All models start from random weights.

The loader memory-maps token files and samples uniformly over valid within-piece window starts. A piece with $L>1{,}024$ tokens provides $L-1{,}024$ starts, so longer pieces receive more sampling opportunities. Inputs are $t_0,\ldots,t_{1023}$ and targets are $t_1,\ldots,t_{1024}$. Neither crosses piece boundaries; shorter pieces cannot supply full training windows.

The loader validates offsets, lengths, and token IDs and overrides the inherited vocabulary default with 376. Architecture configurations and environment details are recorded in the repository.

\begin{table}[H]
\centering
\caption{Shared training settings; the baseline uses dropout 0.1.}
\begin{tabular}{ll@{\qquad}ll}
\toprule Optimizer & AdamW & Weight decay & 0.1\\
$(\beta_1,\beta_2)$ & $(0.9,0.95)$ & Gradient clipping & Norm 1.0\\
Peak / minimum LR & $3\!\times\!10^{-4}$ / $3\!\times\!10^{-5}$ & Seed & 1337\\
Precision & BF16 autocast & Effective batch & 32 sequences\\
Context & 1,024 tokens & Targets/update & 32,768\\\bottomrule
\end{tabular}
\end{table}

AdamW adapts updates using past gradients and separately applies weight decay \citep{loshchilov2019adamw}. Matrix weights receive decay; one-dimensional normalization parameters do not. CUDA uses fused AdamW and, when available, PyTorch scaled dot-product attention. Compilation is disabled and compiler modes were not benchmarked.

The 7,200-second loop budget includes validation and checkpoint writes. Setup and final evaluation/saving add overhead, and an ongoing update finishes at the boundary. This is approximately two hours of loop time, not exactly billed time or fixed FLOPs. Atomic checkpoints save weights, optimizer/scaler, random states, token counts, and elapsed budget. Resuming keeps the original budget; Modal Volumes persist outputs after client disconnection. The image pins Python 3.11, PyTorch 2.6.0, and NumPy 2.2.3.
